% !PRINTED 1
\weberpage{1}
\begin{center}
{\large\gesperrt{INTRODUCTION.}}

\rule{0.18\linewidth}{0.4pt}
\end{center}

We presuppose, in our considerations, the \gesperrt{natural numbers}
$1,2,3,\ldots$ and the rules according to which calculations are made
with these numbers as known and given. The fundamental arithmetical
operations, the so-called \gesperrt{four species}, are \gesperrt{addition},
\gesperrt{multiplication}, to which, as repetition, \gesperrt{exponentiation}
belongs, \gesperrt{subtraction}, and \gesperrt{division}. The first two are
called the \gesperrt{direct operations of calculation}; they are
distinguished by the fact that in the realm of the natural numbers they
can be carried out without restriction. The first of the indirect or
inverse operations, subtraction, can be carried out only when the
minuend is greater than the subtrahend.

The problem of division can be understood in two ways. In the first
elementary interpretation one asks how often the divisor is contained
in the dividend. A number is not contained in a smaller one. But if the
dividend is equal to or greater than the divisor, then the answer to the
question gives a \gesperrt{quotient} and, in most cases, a remainder
which is smaller than the divisor. If no remainder is left, one says
that the division comes out exactly, or that the \gesperrt{dividend} is
divisible by the \gesperrt{divisor}, or that the divisor is a
\gesperrt{factor} or \gesperrt{divisor} of the number which forms the
dividend.

This problem, which already presents itself at the first stages of the
art of calculation, leads to a more deeply lying distinction among
numbers, which is the foundation of all number theory.

Since a factor can never be greater than the number of which it is a
factor, every number has only a finite number of factors. Every number
is divisible by $1$ and by itself, and a number which has no other
divisor is called a \gesperrt{prime-}
%CONT>>
